\chapter{Fonctionnalites Multi-sport}

\section{Concept et architecture}

La plateforme \textit{Sport Academy Pro} est concue pour etre multi-sport, permettant a une seule application de gerer differents sports (football, basketball, handball, speedball, etc.) sans duplication de code.

\section{Principes de conception}

\subsection*{Noyau commun}

Un noyau commun gere les fonctionnalites transversales:

\begin{itemize}[leftmargin=*]
  \item Authentification et gestion des utilisateurs
  \item Systeme de roles et permissions
  \item Calendrier et planning
  \item Messagerie et notifications
  \item Paiements et cotisations
  \item Back-office administratif
\end{itemize}

\subsection*{Elements configurables}

Chaque sport dispose de configurations specifiques:

\begin{itemize}[leftmargin=*]
  \item Positions et roles sur le terrain
  \item Statistiques et indicateurs de performance
  \item Regles et configurations specifiques
  \item Terminologie et labels
\end{itemize}

\section{Implementation technique}

\subsection*{Modele de donnees multi-sport}

\begin{lstlisting}[style=codeblock]
// Entite Sport
@Entity
public class Sport {
    @Id
    @GeneratedValue(strategy = GenerationType.IDENTITY)
    private Long id;

    private String name;           // "Football", "Basketball", etc.
    private String code;           // "FB", "BK", etc.
    private boolean isActive;

    @OneToMany(mappedBy = "sport")
    private List<Position> positions;

    @OneToMany(mappedBy = "sport")
    private List<StatDefinition> statDefinitions;
}

// Entite Position (configurable par sport)
@Entity
public class Position {
    @Id
    @GeneratedValue(strategy = GenerationType.IDENTITY)
    private Long id;

    private String name;           // "Gardien", "Attaquant", etc.
    private String code;           // "GK", "FW", etc.
    private String abbreviation;

    @ManyToOne
    private Sport sport;
}

// Definition des statistiques par sport
@Entity
public class StatDefinition {
    @Id
    @GeneratedValue(strategy = GenerationType.IDENTITY)
    private Long id;

    private String name;           // "Buts", "Passes decisives"
    private String code;           // "goals", "assists"
    private String dataType;       // "integer", "decimal", "time"

    @ManyToOne
    private Sport sport;
}
\end{lstlisting}

\subsection*{Configuration dynamique}

Les configurations sont chargees dynamiquement selon le sport actif:

\begin{lstlisting}[style=codeblock]
@Service
public class SportConfigService {

    public List<Position> getPositions(Long sportId) {
        return positionRepository.findBySportId(sportId);
    }

    public List<StatDefinition> getStatDefinitions(Long sportId) {
        return statDefinitionRepository.findBySportId(sportId);
    }

    public SportConfig getSportConfig(Long sportId) {
        Sport sport = sportRepository.findById(sportId)
            .orElseThrow(() -> new SportNotFoundException(sportId));

        return SportConfig.builder()
            .sport(sport)
            .positions(getPositions(sportId))
            .statDefinitions(getStatDefinitions(sportId))
            .build();
    }
}
\end{lstlisting}

\section{Fonctionnalites multi-sport}

\subsection*{MS01 - Activation ou ajout d'un sport}

L'administrateur peut activer ou desactiver des sports via le back-office:

\begin{itemize}[leftmargin=*]
  \item Liste des sports disponibles
  \item Activation/desactivation en un clic
  \item Configuration initiale des positions et statistiques
\end{itemize}

\subsection*{MS02 - Configuration par sport}

Chaque sport peut etre configure independamment:

\begin{itemize}[leftmargin=*]
  \item Positions specifiques (ex: gardien, defenseur pour le football)
  \item Statistiques pertinentes (buts pour football, points pour basketball)
  \item Regles de calcul et indicateurs
  \item Terminologie et labels
\end{itemize}

\subsection*{MS03 - UX unifiee}

L'interface utilisateur reste la meme quel que soit le sport:

\begin{itemize}[leftmargin=*]
  \item Meme navigation et structure
  \item Seules les donnees et labels changent
  \item Adaptation automatique des formulaires
  \item Filtres et recherches contextuels
\end{itemize}

\section{Exemples d'implementation}

\subsection*{Football}

\begin{table}[H]
\centering
\begin{tabular}{p{0.3\textwidth} p{0.6\textwidth}}
\toprule
\textbf{Element} & \textbf{Configuration} \\
\midrule
Positions & Gardien, Defenseur, Milieu, Attaquant \\
Statistiques & Buts, Passes decisives, Tacles, Interceptions \\
Indicateurs & Note moyenne, Minutes jouees, Cartons \\
\bottomrule
\end{tabular}
\caption{Configuration Football}
\end{table}

\subsection*{Basketball}

\begin{table}[H]
\centering
\begin{tabular}{p{0.3\textwidth} p{0.6\textwidth}}
\toprule
\textbf{Element} & \textbf{Configuration} \\
\midrule
Positions & Meneur, Arriere, Ailier, Pivot \\
Statistiques & Points, Rebonds, Passes decisives, Interceptions \\
Indicateurs & Pourcentage de tir, Plus/minus, Temps de jeu \\
\bottomrule
\end{tabular}
\caption{Configuration Basketball}
\end{table}

\section{Avantages de l'approche multi-sport}

\begin{itemize}[leftmargin=*]
  \item \textbf{Maintenance simplifiee}: Un seul code base a maintenir
  \item \textbf{Extensibilite}: Ajout facile de nouveaux sports
  \item \textbf{Consistance}: UX uniforme pour tous les sports
  \item \textbf{Economie d'echelle}: Developpement et maintenance optimises
\end{itemize}

\section{Integration avec les services IA}

Les services IA s'adaptent automatiquement au sport:

\begin{itemize}[leftmargin=*]
  \item Calcul du potentiel selon les statistiques du sport
  \item Analyse d'evolution contextuelle
  \item Prediction de churn adaptee
  \item Recherche scouting filtree par sport
\end{itemize}